% Part: normal-modal-logic
% Chapter: syntax-and-semantics
% Section: introduction

\documentclass[../../../include/open-logic-section]{subfiles}

\begin{document}

\olfileid[zh]{nml}{syn}{int}

\olsection{引言}

模态逻辑处理的是\emph{模态命题}以及它们之间的衍推关系。模态命题的例子如下：
\begin{enumerate}
\item $2+2=4$ 是必然的。
\item 「明天将要下雨」是必然可能的。
\item 若~$!A$ 是必然可能的，则 $!A$ 是可能的。
\end{enumerate}
可能性与必然性并不是仅有的模态：其他一元联结词也被归为模态，例如「应当如此：$!A$」、「将要如此：$!A$」、「\zhFirst{Dana}{达娜} 知道 $!A$」或「\zhFirst{Dana}{达娜} 相信 $!A$」。

模态逻辑首次出现在亚里士多德的\emph{De Interpretatione}（《解释篇》）中：他是第一个注意到必然性蕴含可能性但反之不然的人；第一个注意到可能性与必然性可以相互定义的人；第一个注意到若 $!A \land !B$ 可能为真，则 $!A$ 可能为真且 $!B$ 可能为真，但反之不然的人；也是第一个注意到若 $!A \to !B$ 是必然的，则若 $!A$ 是必然的，$!B$ 亦然的人。

现代模态逻辑的首次尝试是 C.~I.~\zhFirst{Lewis}{刘易斯} 的工作，其成果为 \zhFirst{Lewis}{刘易斯} 与 Langford 合著的\emph{Symbolic Logic}（《符号逻辑》）（1932 年）。\zhFirst{Lewis}{刘易斯} 与 Langford 对用实质条件句来表示蕴含感到不满：$!A \lif !B$ 是对「$!A$ 蕴含 $!B$」的一种拙劣替代。于是他们提议把蕴含刻画为「必然地，若 $!A$ 则 $!B$」，符号化为 $!A \strictif !B$。在试图厘清这些不同性质的过程中，\zhFirst{Lewis}{刘易斯} 辨认出五个不同的模态系统，\Log{S1}, \ldots, \Log{S4}, \Log{S5}，其中最后两个至今仍在使用。

\zhFirst{Lewis}{刘易斯} 与 Langford 的方法是纯粹\emph{句法的}：他们辨认出合理的公理与规则，并考察在这些手段下什么是可证明的。语义方法在很长一段时间里仍无从得手，直到 Rudolf \zhFirst{Carnap}{卡尔纳普} 在\emph{Meaning and Necessity}（《意义与必然性》）（1947 年）中依据\emph{状态描述}的概念作出首次尝试，即（在该状态描述中）为真的一组原子语句的集合。在把真值定义推广到任意语句 $!A$ 之后，\zhFirst{Carnap}{卡尔纳普} 把 $!A$ 定义为：若它在所有状态描述中都为真，则它是\emph{必然真}的。\zhFirst{Carnap}{卡尔纳普} 的方法无法处理\emph{叠加}模态，因为「可能地必然地 \ldots 可能地 $!A$」这种形式的语句总是化约为最内层的模态算子。

模态语义学上的重大突破来自 Saul \zhFirst{Kripke}{克里普克} 的文章《A Completeness Theorem in Modal Logic》（《模态逻辑中的一个完全性定理》）（JSL，1959 年）。\zhFirst{Kripke}{克里普克} 把他的工作建立在 Leibniz 的观念之上：一个语句若在「所有可能世界处」为真，则它是必然真的。不过，这一观念有着与 \zhFirst{Carnap}{卡尔纳普} 相同的缺陷，即一个语句在世界 $w$ 处（或状态描述 $s$ 中）的真假完全不依赖于 $w$。于是 \zhFirst{Kripke}{克里普克} 假设世界由\emph{可及关系} $R$ 相联，并且「必然地 $!A$」形式的语句在世界 $w$ 处为真当且仅当 $!A$ 在所有从 $w$ \emph{可及}的世界 $w'$ 处为真。提供这一方法某种变体的语义被称为 \zhFirst{Kripke}{克里普克} 语义，它使模态逻辑（复数意义上的诸模态逻辑）得以迅猛发展。

当用 \zhFirst{Kripke}{克里普克} 语义来解释时，模态逻辑向我们展示\emph{关系结构}「从内部看」是什么样子。一个关系结构不过是一个配有一个二元关系的集合（例如，按社会保险号排序的班里学生集合就是一个关系结构）。但事实上，关系结构出现在各种各样的论域中：除了世界诸状态的相对可能性之外，我们还可以有某个主体的认知状态由认知可能性相联，或者一个动力系统的诸状态连同其状态转移，等等。模态逻辑可以用来为所有这些建模：前者给出普通的、真性的模态逻辑；后者给出认知逻辑、动态逻辑等。

我们关注一个特殊的视角，模态逻辑学家称之为「对应理论」。\zhFirst{Kripke}{克里普克} 的早期发现中意义最重大的之一在于：可及关系~$R$ 的许多性质（如它是否传递、是否对称，等等）都可以在\emph{模态语言}之内、借助适当的「模态模式」加以刻画。例如，模态逻辑学家会说，$R$ 的自反性「对应于」模式「若必然地 $!A$，则 $!A$」。我们主要探讨若干由模式 \Ax{D}、\Ax{T}、\Ax{B}、\Ax{4} 与~\Ax{5} 组合而来的经典模态逻辑系统（如 \Log{S4} 与 \Log{S5}）的对应理论。

\end{document}
